\documentclass[aspectratio=169, 11pt]{beamer}
\usepackage{beamer_preamble}

\begin{document}

% ==== Title ===============================================================
\titleslide{Lesson 6-3 \textperiodcentered\ Single-Period Quick}{Logarithms + DOK-3 Spine: q15 (exp\(\leftrightarrow\)log inverse via graph)}

% ==== Do Now ==============================================================
\begin{frame}{Do Now --- Error Analysis + Log Vocab}{Algebra 2 \textperiodcentered\ Lesson 6-3}
\phasebadges{1--2}{7}{bridge to logarithms launch}
\vspace{0.2cm}
\begin{projcard}{Three items --- work solo}{slidegold}
\begin{enumerate}
  \setlength{\itemsep}{0.1cm}
  \item \textbf{q14 Error Analysis:} Find and correct the error in this work:
        \(\displaystyle 16e^t = 98 \;\rightarrow\; e^t = 6.125 \;\rightarrow\; 6.125t = \ln e\)
  \item \textbf{q16 Generalize:} For what values of \(x\) is \(\log_4 x < 0\) true?
  \item \textbf{q3 Vocabulary:} Explain the difference between the common log and the natural log.
\end{enumerate}
\end{projcard}
\end{frame}

% ==== Launch ==============================================================
\begin{frame}{Launch --- Example 3: Evaluate Logarithms}{Algebra 2 \textperiodcentered\ Lesson 6-3}
\phasebadges{1--2}{8}{teacher think-aloud Example 3 (Evaluate Logarithms)}
\vspace{0.15cm}
\begin{projcard}{LOG \(\leftrightarrow\) EXP RULES (keep printed on your page)}{slidegreen}
\(\log_b x = y \;\Leftrightarrow\; b^y = x\) \hfill
COMMON LOG: \(\log x = \log_{10} x\) \hfill
NATURAL LOG: \(\ln x = \log_e x\)

\smallskip
KEY: \(\ln(e^t) = t\) \hspace{1em}
UNDEFINED: \(\log_b 0\) and \(\log_b(\text{negative})\) \hspace{1em}
IDENTITY: \(\log_b(b^k) = k\)
\end{projcard}

\vspace{0.1cm}
\begin{projcard}{Teacher think-aloud --- four cases}{slidegreen}
\small
\begin{tabular}{ll}
\(\log_5 125 = 3\) & \small(\(5^3 = 125\)) \\[2pt]
\(\log_{1/4} 16 = -2\) & \small(\((1/4)^{-2} = 4^2 = 16\); flag: negative exponent) \\[2pt]
\(\log_3 0\): undefined & \small(no power of 3 reaches 0) \\[2pt]
\(\log_2 2^8 = 8\) & \small(identity: \(\log_b(b^k)=k\)) \\
\end{tabular}

\smallskip
\textcolor{slideaccent}{\textbf{30-sec pause:} Rewrite \(\log_4 64\) using the definition with your partner.}
\end{projcard}
\end{frame}

% ==== Explore =============================================================
\begin{frame}{Explore --- TPS + DOK-3 Spine (q15)}{Algebra 2 \textperiodcentered\ Lesson 6-3}
\phasebadges{2--3}{25}{TPS + DOK-3 driver}
\small\textcolor{slidegray}{5 items in order. Plan \(\sim\)10 min for q15 (DOK-3) + 5 min pair presentations.}

\vspace{0.1cm}
\begin{projcard}{Practice \#15 \textperiodcentered\ DOK-3 SPINE (\(\bigstar\))}{slideaccent}
Use the graph of \(y = 3^x\) (\(x \in [-1, 4.5]\), \(y \in [0, 85]\)) to \textbf{estimate} \(\log_3 50\).
Explain your reasoning. \hfill \small\textit{(graph on student packet)}
\end{projcard}
\vspace{-0.12cm}
\begin{projcard}{Practice \#52 \textperiodcentered\ Continuous growth (q52)}{slideblue}
How long for \$250 to grow to \$600 at 4\% APR compounded continuously? Round to nearest year.
\end{projcard}
\vspace{-0.12cm}
\begin{projcard}{Practice \#53 \textperiodcentered\ Compare accounts (q53)}{slideblue}
Michael vs.\ Peter: who reaches \$1{,}800 first? (45-min variant: skip)
\end{projcard}
\vspace{-0.12cm}
\begin{projcard}{Practice \#54 \textperiodcentered\ Richter model (q54)}{slideblue}
\(R = 0.67\log(0.37E)+1.46\). Parts a, b, c. (45-min variant: skip)
\end{projcard}
\vspace{-0.12cm}
\begin{projcard}{Practice \#17 \textperiodcentered\ Proof-style critique (q17)}{slideblue}
Is \(\log_3\!\left(\tfrac{1}{27}\right) = -3\) correct? Use the definition to justify.
\end{projcard}
\end{frame}

% ==== Share / Summary =====================================================
\begin{frame}{Share / Summary --- DOK-3 Reveal}{Algebra 2 \textperiodcentered\ Lesson 6-3}
\phasebadges{2}{5}{academic conversation + 6-4 bridge}
\vspace{0.2cm}
\begin{projcard}{Sentence frame \textperiodcentered\ q15 pair-share}{slideblue}
``On the graph of \(y = 3^x\), I found \(\log_3 50\) by looking for \(y = \underline{\hspace{0.5cm}}\) and reading the x-value.
Since \(3^3 = \underline{\hspace{0.4cm}}\) and \(3^4 = \underline{\hspace{0.4cm}}\), the answer is between \underline{\hspace{0.4cm}} and \underline{\hspace{0.4cm}}, closer to \underline{\hspace{0.6cm}}.
This works because \(\log_3 50 = x\) means \(3^x = \underline{\hspace{0.4cm}}\), so logs and exponentials are \underline{\hspace{0.9cm}}.''
\end{projcard}

\vspace{0.15cm}
\begin{projcard}{Bridge to Lesson 6-4}{slideaccent}
Next lesson: logarithmic \textbf{FUNCTIONS} --- graphing \(y = \log_b x\), domain, range, asymptotes.
Today's log \(\leftrightarrow\) exp definition is the foundation.
\end{projcard}
\end{frame}

% ==== Exit Ticket =========================================================
\begin{frame}{Exit Ticket --- Summary Recap}{Algebra 2 \textperiodcentered\ Lesson 6-3}
\phasebadges{1--2}{5}{summary recap (not graded)}
\vspace{0.2cm}
\begin{projcard}{Complete on your exit ticket}{slidegold}
\begin{enumerate}
  \setlength{\itemsep}{0.25cm}
  \item \(\log_b x = y\) means \underline{\hspace{4cm}}.
  \item \(\log_2 32 = \) \underline{\hspace{1cm}} because \(2^{?} = 32\).
  \item To solve \(e^t = 6.125\), take the \underline{\hspace{1.8cm}} of both sides giving \(t = \) \underline{\hspace{1.6cm}}.
  \item One biggest thing I learned today: \underline{\hspace{5cm}}
\end{enumerate}
\end{projcard}
\end{frame}

\end{document}
